\begin{textsample}{POS Dim 4 – ba – Score 18.00 – ba/ba\_em\_7.txt}  \label{ex:f4_pos_023}
5.2.
EDUCAÇÃO NA DIVERSIDADE E PARA A DIVERSIDADE O cenário educacional brasileiro perpassa por significativas mudanças no que se refere às políticas públicas.
A diversidade na educação é uma demanda importante e que se desponta no cenário das políticas públicas educacionais, sobretudo, para se fazer cumprir o direito à educação de qualidade social para todas as pessoas.
Uma dessas políticas educacionais estruturantes, é o currículo, pois, a partir dele, são declarados os valores e conhecimentos que determinado sistema, rede ou instituição de ensino pretende conduzir uma proposta de formação de pessoas.
Mas, para considerar a diversidade no contexto do “ currículo vivo ” das escolas, faz -se necessário compreender diversidade, que, segundo Nilma Lino Gomes : > a diversidade pode ser entendida como a construção histórica, cultural e social das diferenças.
A construção das diferenças ultrapassa as características biológicas, observáveis a olho nu.
As diferenças são também construídas pelos sujeitos sociais ao longo do processo histórico e cultural, nos processos de adaptação do \textbf{homem} e da \textbf{mulher} ao meio social e no contexto das relações de poder.
Sendo assim, mesmo os aspectos tipicamente observáveis, que aprendemos a ver como diferentes desde o nosso nascimento, só passaram a ser percebidos dessa forma, porque nós, seres humanos e sujeitos sociais, no contexto da cultura, assim os nomeamos e identificamos. ( GOMES, 2008 ).
A escola é, em sua essência, um local diverso, de corpos, de \textbf{identidades}, de pertencimentos étnico-raciais de gerações, de saberes singulares e diversos.
Tal ambiente envolve uma variedade cultural apresentada não apenas por diversos formatos linguísticos, mas também por diversas demonstrações de tradições, crenças, culturas, valores, posicionamentos políticos, expressões, orientações afetivo-sexuais, \textbf{identidades} de gênero, sexualidades e etnias variadas.
Uma educação que versa acerca da diversidade tem como base, ainda, o respeito às suas municipalidades expressas no cotidiano escolar.
Sendo assim, é muito importante que os/as educadores/as busquem por práticas educativas que valorizem e promovam a cultura e os aspectos indenitários dos municípios que integram os \textbf{Territórios} de \textbf{Identidade} do \textbf{estado}, “ suas características geoambientais e socioambientais, bem como com a sociedade, a história, a cultura, a economia ” ( BAHIA, 2019 ).
Dessa forma, a escola preservará as “ histórias, memórias e \textbf{identidades} da história local nas diferentes temporalidades ” ( BAHIA, 2019 ) e tais memórias e \textbf{identidades} locais podem ser contextualizadas para aprendizagem de outros conteúdos curriculares.
As diretrizes apresentadas neste eixo mostram que as políticas públicas implementadas no \textbf{Estado} da Bahia podem impulsionar e promover a construção de uma escola mais plural.
A implementação, na rede estadual, de elementos como rodas discursivas e demais práticas pedagógicas que sejam inclusivas vão auxiliar os profissionais de educação que atuam na educação básica – seja em escolas públicas ou privadas do \textbf{estado} – a compreender que incluir pessoas diferentes no cotidiano escolar é, também, um ato de respeito.
Contribuir com e para a aceitação do/a aluno/a, de uma forma que seja construtiva para estes/estas, tem como objetivo, ainda, fazer valer o que descreve a Constituição de 1988 no Art. 3º, que discorre sobre os objetivos fundamentais da República Federativa do Brasil : I.
Construir uma sociedade livre, justa e solidária ; II.
Garantir o desenvolvimento nacional ; III.
Erradicar a pobreza e a marginalização e reduzir as desigualdades sociais e regionais ; IV.
Promover o bem de todos, sem preconceitos de origem, \textbf{raça}, sexo, cor, idade quaisquer outras formas de \textbf{discriminação}. ( BRASIL, 1988 ).
O debate de questões que envolvem a diversidade e os direitos humanos, como as relações étnico-raciais, as diversidades \textbf{religiosas} - em especial, as oriundas de matriz africanas e afro-brasileiras -, os povos \textbf{indígenas}, os povos \textbf{ciganos}, as pessoas com deficiências ( PCD ), as \textbf{identidades} de gênero e sexualidades, bem como o enfrentamento a todas as formas de violência e de \textbf{discriminação} a todos os grupos sociais, considerados minoritários, devem ser pauta ou abordados de forma contextualizadas na elaborações de atividades pedagógicas que \textbf{atravessam}, intersecionam, promovam e possibilitam novas formas de pensar e vivenciar a educação para diversidade nas instituições de ensino, de forma leve e que não provoque sofrimento a esses grupos.
As dimensões supramencionadas, que tratam da diversidade humana, presentes em todos os espaços sociais, são contemplados pelo conceito de supralegalidade do ordenamento jurídico brasileiro, conforme ressalta a estratégia 2.12 do Plano Estadual de Educação ( PEE ) – Lei nº 13.559/2016 : > estimular que o respeito às diversidades seja objeto de tratamento transversal pelos professores, bem como pelas Instituições de Ensino Superior nos currículos de graduação, respeitando os Direitos Humanos e o combate a todas as formas de \textbf{discriminação} e \textbf{intolerância}, à luz do conceito de supralegalidade presente no ordenamento jurídico brasileiro ( Bahia, 2016 ).
No campo da diversidade de gênero e sexualidades, outro normativo recente que reforça a atuação dos/as educadores/as sobre a Educação das Relações de Gêneros e Sexualidades, é a publicação da Resolução CEE-BA nº 45/2020, que dispõe sobre Educação das Relações de Gêneros e Sexualidades, a ser observada pelas instituições públicas e privadas integrantes do Sistema Estadual de Ensino da Bahia, conforme Art. 2º da referida Resolução que tem como princípios : I.
Dignidade humana ; II.
Laicidade da educação ; III.
Educação democrática ; IV.
A equidade de gênero ; V.
Reconhecimento e valorização das diferenças e da diversidade de gênero e sexualidade ; VI.
Pluralismo de idéias e concepções pedagógicas ; VII.
Transversalidade e contextualização do conhecimento.
Neste eixo estruturante do DCRB, a centralidade do currículo na educação para as relações étnico-raciais é reforçada a partir da aplicabilidade da Lei nº 10.639/2003, que trata da inserção da História e Cultura Africana e Afro-brasileira nos currículos escolares, bem como da aplicabilidade da Lei nº 11.645/08, que insere os estudos da História e Cultura \textbf{Indígena} nos currículos das escolas.
A importância dos povos africanos e \textbf{indígenas} deve constituir a “ espinha dorsal ” para a organização curricular e orientar as práticas pedagógicas das escolas do \textbf{Estado}.
Faz -se necessário considerar a história da \textbf{ancestralidade} africana e \textbf{indígena}, evidenciando as contribuições desses povos nos mais diversos campos do conhecimento e do saber, suas estratégias de resistência e todas as conquistas realizadas, ao longo do processo de construção do país. > o \textbf{Estado} adotará ações para assegurar a qualidade do ensino da História e da Cultura Africana, Afro-brasileira e \textbf{Indígena} nas unidades do Ensino Fundamental e Médio do Sistema Estadual de Ensino, em conformidade com o estabelecido pela Lei de Diretrizes e Bases da Educação Nacional, assegurando a estrutura e os meios necessários à sua efetivação, inclusive no que se refere à formação permanente de educadores, realização de campanhas e disponibilização de material didático específico, no contexto de um conjunto de ações integradas com o combate ao \textbf{racismo} e à \textbf{discriminação} \textbf{racial} nas escolas ( BAHIA, 2014 ).
Um outro aspecto importante dentro deste eixo é promover diálogos com as multiplicidades das denominações \textbf{religiosas}, presentes nas escolas, considerando que o \textbf{estado} brasileiro é \textbf{laico}, objetivando a convivência respeitosa entre as pessoas, conforme dispõe o inciso VI do Art. 5º da Constituição que diz : “ é inviolável a liberdade de consciência e de crença, sendo assegurado o livre exercício dos cultos \textbf{religiosos} e garantida, na forma da lei, a proteção aos locais de culto e a suas liturgias ( BRASIL, 1988 ).
Também, Destarte a UNESCO se refere a diversidade cultural como necessária para a conceituação do gênero humano como a diversidade biológica o é para a natureza.
Constituindo, assim, um patrimônio comum da humanidade e deve ser reconhecida e consolidada em benefício das gerações presentes e futuras.
Deste modo, os/as alunos/as são a expressão do futuro que queremos e, para isso, é necessário que, no presente, trabalhemos uma educação na e para a diversidade, valorizando e respeitando a integridade do singular humano nas mais distintas formas de expressões que possam surgir e, mesmo que não saibamos “ dar ” nome a “ estas ” determinadas situações, que possamos estar abertos e abertas a aprender e apreender a fazer um espaço educativo melhor.
Assim, a educação na e para a diversidade deve vislumbrar um lugar em meio à dinâmica social, promovendo trocas de experiências, diálogos, rodas de conversas, planejamentos pedagógicos, a partir das realidades existentes no “ \textbf{chão} ” da escola, envolvendo os/as atores/atrizes da educação e toda esta nova emergência dos conceitos sobre : as pluralidades, multiplicidades, variedades, diferenças e heterogeneidade da tradição que contrapõem -se aos novos modelos apresentados nas diversas formas de interpretar a nossa existência.

% matched lemmas: ancestralidade, atravessar, chão, cigano, discriminação, estado, homem, identidade, indígena, intolerância, laico, mulher, racial, racismo, raça, religioso, território
\end{textsample}
